\documentclass[12pt]{article}
\input{iati_structure.tex}
\renewcommand{\bottomtitlespace}{2cm}

\usepackage[english]{babel}

\begin{document}

\renewcommand{\headerLang}{Russian}
\renewcommand{\problemName}{AB23. ТРЕУГОЛЬНИК}

\renewcommand{\headerLeft}{IATI Day 2, Senior group}
\renewcommand{\headerRightFirst}{XVII INTERNATIONAL ADVANCED TOURNAMENT IN INFORMATICS}
\renewcommand{\headerRightSecond}{ONLINE 2026}

\renewcommand{\tl}{$10$ сек.}
\renewcommand{\ml}{$1024$ МБ}
\problem{Задача \problemName}
\addauthor{Автор: Борис Михов}{2026}{04}{19}

Новое руководство национального комитета в лице координаторов групп A и B желает придерживаться строгой программы по задачам для отбора национальной команды. Для этой цели им потребуется дать геометрическую задачу. Более того — количество интерактивных задач, данных на текущий момент, ниже стандарта. Чтобы исправить кризис в отборе на \texttt{IOI}, Сашка решил дать задачу, которая одновременно является и интерактивной, и геометрической. Её условие следующее:

Жюри загадало перестановку $P_0,P_1,...,P_{N-1}$ чисел $\{1,2,3,...,N\}$. Вы должны найти эту перестановку. Для этого вы можете задать жюри следующий вопрос: «Можно ли составить треугольник положительной площади со сторонами $P_A,P_B,P_C$?»

Заметим, что (по неравенству треугольника) это возможно тогда и только тогда, когда:

\[P_A + P_B > P_C,\ P_A + P_C > P_B \quad \text{и} \quad P_B + P_C > P_A.\]

Напишите программу \textbf{\texttt{triangle}}, содержащую функцию \texttt{solve}, которая будет компилироваться вместе с программой жюри и взаимодействовать с ней, задавая вопросы описанного выше типа. По завершении своей работы она должна определить перестановку.

\subsection{Детали реализации}
Вы должны реализовать функцию \texttt{solve}:
\begin{lstlisting}
 std::vector<int> solve(int N)
\end{lstlisting}
\begin{itemize}
	\item $N$: длина перестановки.
\end{itemize}

Эта функция будет вызываться $T$ раз на тест — по одному разу на подтест, каждый раз с одинаковым $N$, и должна возвращать загаданную перестановку для подтеста. Для этого ваша программа может вызывать функцию жюри \texttt{query}:
\begin{lstlisting}
 bool query(int A, int B, int C)
\end{lstlisting}
\begin{itemize}
	\item $A$, $B$, $C$: индексы сторон $P_A, P_B, P_C$.
\end{itemize}

Функция возвращает \texttt{true}, если существует треугольник положительной площади со сторонами $P_A,P_B,P_C$, и \texttt{false} в противном случае.

\newpage

\subsection{Ограничения}
\vspace{0.1em}
\begin{itemize}
	\item $N = T = 1~000$
\end{itemize}

\subsection{Подзадачи}
\begin{table}[H]
	\begin{tblr}{width=\linewidth,colspec={|Q[c,m]|Q[c,m]|X[c,m]|}}
		\hline
		\textbf{Подзадача} & \textbf{Баллы} & \textbf{Описание}\\
		\hline
		$1$ & $100$ & Подзадача состоит из $1$ теста с $T=1~000$ случайно сгенерированными перестановками, каждая из $N=1~000$ элементов.  \\ 
		\hline
	\end{tblr}
	\caption*{Баллы за данную подзадачу начисляются только при прохождении всех её тестов.}
\end{table}
\FloatBarrier

\subsection{Оценка}

Пусть $Q_{contestant}$ — среднее количество запросов, которое ваша программа делает за один вызов \texttt{solve} на единственном тесте, и пусть $Q_{target}=8770$. Тогда ваша оценка за задачу будет:

\[
\begin{cases}
0 & \text{если } Q_{contestant} > 2\times 10^6, \\
100 & \text{если } Q_{contestant} \le Q_{target}, \\
100 \times \max\left(0.15,\; 1-\sqrt{1-\left(\frac{Q_{target}}{Q_{contestant}}\right)^{0.55}}\right) & \text{если } Q_{contestant} > Q_{target}.
\end{cases}
\]

\newpage

\subsection{Пример взаимодействия}
Для этого примера взаимодействия $N=4$, $T=2$.
\begin{table}[H]
	\begin{tblr}{|c|c|}
		\hline
		\textbf{Действие участника} & \textbf{Действие жюри} \\
		\hline
		& \texttt{solve(4)} \\
		\hline
		\texttt{query(0, 0, 0)} & \texttt{true} \\
		\hline
        \texttt{query(0, 1, 2)} & \texttt{false} \\
		\hline
        \texttt{query(0, 1, 3)} & \texttt{false} \\
		\hline
        \texttt{query(0, 2, 3)} & \texttt{true} \\
		\hline
        \texttt{query(1, 2, 3)} & \texttt{false} \\
		\hline
        \texttt{return \{3, 1, 2, 4\};} & \\
        \hline
        & \texttt{solve(4)} \\
		\hline
		\texttt{query(0, 1, 2)} & \texttt{false} \\
		\hline
        \texttt{query(0, 1, 3)} & \texttt{true} \\
		\hline
        \texttt{query(0, 2, 3)} & \texttt{false} \\
		\hline
        \texttt{query(1, 2, 3)} & \texttt{false} \\
		\hline
        \texttt{return \{4, 2, 1, 3\};} & \\
		\hline
	\end{tblr}
\end{table}
\FloatBarrier
	
\subsection{Локальное тестирование}
Формат ввода:
\begin{itemize}
	\item строка $1$: три целых числа $T$, $N$ и $R$ — размер перестановок, количество подтестов и режим выполнения. Если $R = 1$, локальный проверяющий сгенерирует перестановки равномерно случайно и ожидает число на первой строке — $S$, которое будет seed'ом для его генератора случайных чисел. Если $R = 2$, ввод продолжается следующим образом:
	\item строки $2$ по $(1+T)$: перестановки чисел $\{1,2,\dots,N\}$.
\end{itemize}

Формат вывода:
\begin{itemize}
	\item строка $1$: сообщение об ошибке или среднее количество запросов по подтестам, если все перестановки определены верно.
\end{itemize}

\end{document}